% ID: SEL-D0062-OB-E02
% ESTADO: cotejada. Segunda pasada visual y borrador compilado revisados.
% Fuente: Selectividad/sel_2010_matematicas_aplicadas/mat_apli/mat_apli_junio/mat_apli_junio_ra/EXAMEN-ANDALUCÍA-4.pdf
% SHA-256: 4e10e40d60aebdab7570fa77b9ce953d7bf6c9a4c8869a5d3f9b7ca3c907c154
% Página: 2; opción B; ejercicio original 2.

Un depósito lleno de agua se vacía por un sumidero que tiene en la parte baja. El volumen de agua, en $m^3$, que hay en cada momento en el depósito, desde que empieza a vaciarse, viene dado por la función $V(t)=8-t+\dfrac{t^2}{32}$, donde $t$ es el tiempo en minutos.
\begin{enumerate}[label=\alph*)]
\item \textbf{(0.5 puntos)} \textquestiondown{}Cuál es la capacidad del depósito?
\item \textbf{(0.5 puntos)} \textquestiondown{}Cuánto tiempo tarda en vaciarse?
\item \textbf{(0.8 puntos)} Represente gráficamente la función $V$.
\item \textbf{(0.7 puntos)} Calcule la derivada de esa función en $t=8$ e interprete su significado.
\end{enumerate}
